\documentclass[11pt]{article}
\usepackage{amsmath}
\usepackage{amssymb}
\usepackage{geometry}
\usepackage{array}
\usepackage{longtable}
\geometry{margin=1in}

\setcounter{secnumdepth}{0}
\newcommand{\cmark}{$\checkmark$}
\newcommand{\xmark}{$\times$}

\title{AMC 12 Teacher Report: April 23--26, 2026}
\author{}
\date{}

\begin{document}
\maketitle

\section{Grading Summary}

\textbf{Scope.} Rows 1--15 of \texttt{amc12-2026-04\_student.csv} were completed. Rows 16--50 were not attempted yet and were left blank in the redo CSV.

\textbf{Timing standard.} Use the student's AMC12 pacing rule: Questions 1--10 should average 2:00 each, Questions 11--15 should average 3:00 each, Questions 16--20 should stay near a 4:00 cap, and selected Questions 21--25 may receive up to 5:00 only when progress is clear. Since this worksheet contains Problem 10 items, the target budget is 2:00 per problem. Correct answers over 2:00 are treated as speed-review items, not redo items.

\section{AMC12 Pace Strategy}

\begin{longtable}{@{}>{\raggedright\arraybackslash}p{0.24\textwidth} >{\centering\arraybackslash}p{0.18\textwidth} >{\centering\arraybackslash}p{0.18\textwidth} >{\raggedright\arraybackslash}p{0.32\textwidth}@{}}
\textbf{Question Band} & \textbf{Target} & \textbf{Checkpoint} & \textbf{Decision Rule}\\
\hline
Q1--Q10 & 2:00 each & 20:00 total & Build a time bank. If there is no clear path after 90 seconds, mark it and move.\\
Q11--Q15 & 3:00 each & 35:00 total & Solve cleanly, but skip if still setting up after 2.5 minutes.\\
Q16--Q20 & Up to 4:00 each & 55:00 total & Attempt only while progress is visible; otherwise save the time for later.\\
Q21--Q25 & Selective, up to 5:00 & Final 20:00 & Do not spend 5 minutes on all five. Choose accessible problems and return to flagged work.\\
\end{longtable}

\textbf{Important caveat.} A flat allocation of 2 minutes for Q1--Q10, 3 minutes for Q11--Q15, and 5 minutes for every Q16--Q25 problem totals \(20+15+50=85\) minutes, which exceeds the 75-minute AMC12 limit. Therefore the 5-minute rule for Q16--Q25 is a maximum for selected problems, not a promise to attempt all ten hard problems fully.

\small
\begin{longtable}{@{}>{\raggedright\arraybackslash}p{0.05\textwidth} >{\raggedright\arraybackslash}p{0.26\textwidth} >{\centering\arraybackslash}p{0.08\textwidth} >{\centering\arraybackslash}p{0.08\textwidth} >{\centering\arraybackslash}p{0.08\textwidth} >{\centering\arraybackslash}p{0.08\textwidth} >{\raggedright\arraybackslash}p{0.17\textwidth}@{}}
\textbf{\#} & \textbf{Source} & \textbf{Stud.} & \textbf{Key} & \textbf{Time} & \textbf{Result} & \textbf{Note}\\
\hline
1 & 2000 AMC 12 P10 & E & E & 03:27 & \cmark & Over budget\\
2 & 2001 AMC 12 P10 & D & D & 01:09 & \cmark & Efficient\\
3 & 2002 AMC 12A P10 & D & D & 01:42 & \cmark & Efficient\\
4 & 2002 AMC 12B P10 & B & A & 01:45 & \xmark & Redo\\
5 & 2003 AMC 12A P10 & D & D & 01:23 & \cmark & Efficient\\
6 & 2003 AMC 12B P10 & B & B & 00:31 & \cmark & Very fast\\
7 & 2004 AMC 12A P10 & C & C & 03:58 & \cmark & Over budget\\
8 & 2004 AMC 12B P10 & A & A & 00:37 & \cmark & Very fast\\
9 & 2005 AMC 12A P10 & B & B & 05:14 & \cmark & Over budget\\
10 & 2005 AMC 12B P10 & E & E & 02:06 & \cmark & Slightly over\\
11 & 2006 AMC 12A P10 & D & E & 00:40 & \xmark & Redo\\
12 & 2006 AMC 12B P10 & A & A & 00:45 & \cmark & Very fast\\
13 & 2007 AMC 12A P10 & A & A & 02:10 & \cmark & Slightly over\\
14 & 2007 AMC 12B P10 & C & C & 03:05 & \cmark & Over budget\\
15 & 2008 AMC 12A P10 & D & D & 01:05 & \cmark & Efficient\\
\end{longtable}
\normalsize

\textbf{Completed score:} 13 correct out of 15 attempted, or \(86.7\%\). \textbf{Redo items:} rows 4 and 11. \textbf{Speed-review items:} rows 1, 7, 9, 10, 13, and 14.

\section{Speed-Review Solutions}

\subsection{Row 1: 2000 AMC 12 Problem 10}

\textbf{Answer:} E.

\textbf{Fast method.} Write each transformation as a coordinate rule instead of trying to visualize the whole 3D motion.
\[
(x,y,z)\xrightarrow{\text{reflect in }xy}(x,y,-z)
\xrightarrow{180^\circ\text{ about }x}(x,-y,-z)
\xrightarrow{+5\text{ in }y}(x,-y+5,-z).
\]
Applying this to \(P=(1,2,3)\) is fast:
\[
(1,2,3)\to(1,2,-3)\to(1,-2,3)\to(1,3,3).
\]

\textbf{Speed focus.} Use a three-column sign table: \(x\) unchanged, \(y\) changes only under the rotation and translation, and \(z\) changes under the reflection and rotation. This should be a 60--90 second problem once the rules are memorized.

\subsection{Row 7: 2004 AMC 12A Problem 10}

\textbf{Answer:} C.

\textbf{Fast method.} For an odd number of consecutive integers, the median equals the average. There are \(49=7^2\) integers and the sum is \(7^5\), so
\[
\text{median}=\frac{7^5}{49}=\frac{7^5}{7^2}=7^3.
\]

\textbf{Speed focus.} The key recognition is ``odd consecutive list \(\Rightarrow\) median = mean.'' Do not find endpoints. Pair terms around the center mentally, then divide sum by count.

\subsection{Row 9: 2005 AMC 12A Problem 10}

\textbf{Answer:} B.

\textbf{Fast method.} Count faces, not small cubes. The total number of faces on all \(n^3\) unit cubes is \(6n^3\). The painted red faces are exactly the surface unit squares of the original cube, so there are \(6n^2\). Thus
\[
\frac{6n^2}{6n^3}=\frac{1}{n}=\frac14,
\]
so \(n=4\).

\textbf{Speed focus.} When a painted-solid problem asks for a fraction of faces, first compare ``surface area faces'' to ``all cube faces.'' Avoid classifying corner, edge, and face-center cubes unless the problem asks how many cubes have 1, 2, or 3 painted faces.

\subsection{Row 10: 2005 AMC 12B Problem 10}

\textbf{Answer:} E.

\textbf{Fast method.} Compute only until a cycle appears:
\[
2005\to 2^3+5^3=133,\qquad
133\to 1^3+3^3+3^3=55,\qquad
55\to 5^3+5^3=250,
\]
and
\[
250\to 2^3+5^3=133.
\]
So from term 2 onward the cycle is \(133,55,250\), with period 3. Term 1 is outside the cycle, so for \(k\ge2\) index into \([133,55,250]\) by \((k-2)\bmod 3\). Since \(2005-2=2003\equiv2\pmod3\), the 2005th term is the third entry of the cycle, \(250\).

\textbf{Speed focus.} For digit-iteration sequences, do not continue linearly. Find the first repeat, write the cycle, then use modular arithmetic.

\subsection{Row 13: 2007 AMC 12A Problem 10}

\textbf{Answer:} A.

\textbf{Fast method.} A \(3:4:5\) triangle is right, so its hypotenuse is the diameter of the circumcircle. The circle radius is 3, so the hypotenuse is 6. The scale factor from \(3:4:5\) is \(6/5\), giving legs
\[
\frac{18}{5}\quad\text{and}\quad \frac{24}{5}.
\]
The area is
\[
\frac12\cdot\frac{18}{5}\cdot\frac{24}{5}=\frac{216}{25}=8.64.
\]

\textbf{Speed focus.} For any right triangle in a circle, immediately use ``hypotenuse = diameter.'' This avoids trigonometry and side-solving.

\subsection{Row 14: 2007 AMC 12B Problem 10}

\textbf{Answer:} C.

\textbf{Fast method.} The total group size does not change: two girls leave and two boys arrive. The fraction of girls drops from \(40\%\) to \(30\%\), a drop of \(10\%\) of the total group. That drop is exactly 2 girls, so
\[
10\%\text{ of total}=2 \quad\Rightarrow\quad \text{total}=20.
\]
Initially \(40\%\) of 20 were girls, so the answer is \(8\).

\textbf{Speed focus.} Compare the before/after percentages before setting up equations. If the total is unchanged, a percentage drop can be converted directly into a number of people.

\section{Redo Solutions}

\subsection{Row 4: 2002 AMC 12B Problem 10}

\textbf{Correct answer:} A. The student chose B.

\textbf{Systematic method.} The set
\[
\{1,4,7,10,13,16,19\}
\]
is an arithmetic sequence. Every element has the form \(1+3k\), where \(k=0,1,2,3,4,5,6\). The sum of three distinct elements is
\[
(1+3a)+(1+3b)+(1+3c)=3+3(a+b+c).
\]
So different original sums correspond exactly to different values of \(a+b+c\), where \(a,b,c\) are distinct numbers chosen from \(0\) through \(6\).

The smallest possible index sum is \(0+1+2=3\), and the largest is \(4+5+6=15\). All values from 3 through 15 are possible, so there are
\[
15-3+1=13
\]
different sums.

\textbf{Methodology lesson.} For arithmetic-sequence choice problems, convert to indices and count the possible index sums. Do not count triples; many different triples can produce the same integer sum. Always separate ``number of selections'' from ``number of distinct values.''

\subsection{Row 11: 2006 AMC 12A Problem 10}

\textbf{Correct answer:} E. The student chose D.

\textbf{Systematic method.} Let
\[
\sqrt{120-\sqrt{x}}=k,
\]
where \(k\) is a nonnegative integer. Then
\[
120-\sqrt{x}=k^2
\quad\Rightarrow\quad
\sqrt{x}=120-k^2.
\]
Since \(\sqrt{x}\ge 0\), we need \(120-k^2\ge 0\), so \(k^2\le 120\). The possible values are
\[
k=0,1,2,\ldots,10,
\]
because \(10^2=100\le120\) and \(11^2=121>120\). This gives 11 possible values of \(k\), and each one produces exactly one real value
\[
x=(120-k^2)^2.
\]

\textbf{Methodology lesson.} For nested radicals, name the whole expression first, then translate every square root into a domain inequality. Include endpoint cases: here \(k=0\) is valid, and missing it changes 11 to 10.

\section{Student Strengths}

\begin{itemize}
  \item Strong accuracy on the attempted set: 13 out of 15 correct on mid-band Problem 10 material.
  \item Very fast execution on several algebra/geometry items: rows 6, 8, 11, and 12 were under one minute. Row 11 was wrong, but the speed shows the student is willing to commit; the next step is adding a quick endpoint check before bubbling.
  \item Good breadth across topics: coordinate geometry, ratios, sequences, counting, geometry, cube face counting, recursive digit process, inequalities, triangle geometry, and rate equations were mostly handled correctly.
\end{itemize}

\section{Growth Opportunities}

\begin{itemize}
  \item Add a 10-second classification step before solving: ``Am I counting choices, values, cases, or arrangements?'' This directly targets row 4.
  \item For radical/counting-domain problems, write the auxiliary variable and endpoint inequalities before answering. This directly targets row 11.
  \item Improve speed on straightforward recognition problems. Rows 1, 7, 9, 10, 13, and 14 were correct but over the 2-minute Q10 pace; each has a one-line shortcut once the problem type is recognized.
  \item Use a final ``endpoint and units'' check. For row 11, check both endpoints \(k=0\) and \(k=10\); an answer of 10 usually means one endpoint was dropped. For row 9, verify the numerator and denominator count the same object, unit-cube faces, before simplifying.
\end{itemize}

\section{Recommended Next Practice}

\begin{itemize}
  \item Redo rows 4 and 11 without looking at the solution, then explain the counting object in one sentence.
  \item Re-solve rows 1, 7, 9, 10, 13, and 14 with a 2:00 timer. The goal is not just to get the answer again, but to use the faster method.
  \item For future Problem 10 sets, pause at 1:30. If the path is not clear, switch to answer-choice testing, invariant recognition, or a simpler variable setup.
\end{itemize}

\end{document}
